% CAPÍTULO 2 -- OBJETIVOS (VERSÃO CONCISA TCC UESC)
% ==========================================================================
\chapter{Objetivos}
\label{cap:objetivos}

\section{Objetivo geral}
\label{sec:obj-geral}

O objetivo geral deste trabalho de conclusão de curso consiste em conceber, modelar, otimizar parametricamente e qualificar estruturalmente um chassi primário para nanossatélite no padrão CubeSat 1U integrado por painéis celulares auxéticos reentrantes em liga de alumínio aeroespacial AlSi10Mg fabricado por Manufatura Aditiva via Fusão em Leito de Pó por Laser (L-PBF), maximizando o alívio de massa e a atenuação elastodinâmica passiva na carga útil sob os perfis estocásticos de vibração aleatória da norma NASA GSFC-STD-7000A.

\section{Objetivos específicos}
\label{sec:obj-especificos}

Para atingir a meta central da pesquisa, estruturam-se nove objetivos específicos distribuídos em quatro fases metodológicas interligadas:

\subsection*{Fase 1: Concepção geométrica e portões de manufaturabilidade aditiva (DfAM)}
\begin{enumerate}
  \item \textbf{Parametrização dimensional da célula unitária auxética reentrante:} formular analítica e computacionalmente a célula bidimensional e tridimensional governada pelas variáveis de projeto $\mathbf{x} = \begin{bmatrix} \theta & t & l & h \end{bmatrix}^{\top}$, mapeando a densidade relativa celular $\rho_{\text{rel}}$ e a cinemática de Poisson negativo ($\nu_{\text{eff}} < 0$).
  \item \textbf{Implementação de portões automáticos de conformidade DfAM:} validar restrições estritas do processo L-PBF em AlSi10Mg, impondo sobre-salha autossuportada ($\theta_{\text{overhang}} \ge \SI{35.0}{\degree}$), espessura mínima de costela para consolidação da poça de fusão ($t \ge \SI{0.50}{\milli\metre}$) e vão livre desobstruído para evacuação de pó não sinterizado ($\text{gap}_{\text{despoeiramento}} \ge \SI{1.50}{\milli\metre}$ e $d_{\text{drain}} \ge \SI{2.00}{\milli\metre}$).
  \item \textbf{Síntese tridimensional paramétrica e exportação neutral B-Rep:} integrar a geometria dos quatro trilhos monolíticos de contato com o dispensador P-POD (\rail{}) e exportar modelos no formato neutro STEP AP214 (ISO 10303-21) via kernel aberto OpenCASCADE.
\end{enumerate}

\subsection*{Fase 2: Modelagem fenomenológica e elastodinâmica}
\begin{enumerate}[resume]
  \item \textbf{Modelagem das propriedades mecânicas efetivas por Gibson-Ashby:} avaliar analiticamente as leis de escala do módulo elástico longitudinal ($E^*$) e limite de escoamento celular ($\sigma_y^*$) para sólidos celulares sob flexão de costela.
  \item \textbf{Análise de atenuação acústica e confinamento de ondas elásticas:} verificar a resposta elastodinâmica passiva e os mecanismos de descasamento de impedância acústica ($Z = \rho c$) nas interfaces entre os painéis auxéticos esbeltos e os trilhos maciços monolíticos.
\end{enumerate}

\subsection*{Fase 3: Esteira computacional aberta de alta fidelidade (CAE)}
\begin{enumerate}[resume]
  \item \textbf{Discretização espacial com elementos finitos de alta ordem (Gmsh Tet10):} automatizar a geração de malhas tetraédricas de 10 nós de segunda ordem com refinamento através da espessura ($h \le t/3$) e nos entalhes reentrantes críticos ($R_{\text{notch}} \ge \SI{0.50}{\milli\metre}$).
  \item \textbf{Orquestração de análises modais e espectrais PSD no Code\_Aster e CalculiX:} implementar a extração de frequências naturais fundamentais ($f_1$) via algoritmo de Sorensen/Lanczos e resolver a resposta à vibração aleatória (\SIrange{20}{2000}{\hertz}, \SI{14.1}{\Grms}) via superposição modal estocástica com amortecimento constante ($\zeta = \SI{2.0}{\percent}$).
\end{enumerate}

\subsection*{Fase 4: Metamodelagem neural, otimização evolutiva e qualificação espacial}
\begin{enumerate}[resume]
  \item \textbf{Desenvolvimento de modelo substituto neural guiado por física (\textit{Physics-Guided ResNet}):} treinar uma rede neural profunda com blocos residuais sobre o banco de dados do planejamento de experimentos (DoE de 250 pontos gerado via Hipercubo Latino e modelo multiescala calibrado), incorporando termos de perda informados por física de Gibson-Ashby e monotonicidade de rigidez.
  \item \textbf{Otimização multiobjetivo (NSGA-II) e qualificação aeroespacial por TOPSIS:} determinar a fronteira de Pareto ótima e selecionar a configuração de voo (\#28) satisfazendo simultaneamente os critérios da NASA: $f_1 \ge \SI{100.0}{\hertz}$, Margem de Segurança positiva ($MS_{\text{yield}} > 0$) e acúmulo de dano por fadiga de Steinberg qualificado ($D \le \num{0.25}$).
\end{enumerate}

\section{Critérios de aceitação e rastreabilidade metodológica}
\label{sec:criterios-aceitacao}

A Tabela~\ref{tab:rastreabilidade-tcc} sintetiza a matriz de rastreabilidade entre as fases metodológicas, os critérios quantitativos de aceitação e as normas espaciais de referência.

\begin{table}[htb]
  \caption{Matriz de rastreabilidade e critérios quantitativos de aceitação do projeto de TCC}
  \label{tab:rastreabilidade-tcc}
  \centering
  \small
  \SingleSpacing
  \begin{tabularx}{\textwidth}{@{}>{\centering\arraybackslash}p{1.4cm} >{\raggedright\arraybackslash}p{4.2cm} >{\raggedright\arraybackslash}p{4.2cm} >{\centering\arraybackslash}X@{}}
    \toprule
    \textbf{Fase} & \textbf{Requisito / Parâmetro} & \textbf{Critério de Aceitação} & \textbf{Norma / Referência} \\
    \midrule
    Fase 1 & Ângulo autossuportado L-PBF    & $\theta_{\text{overhang}} \ge \SI{35.0}{\degree}$ & NASA-STD-6030 \\
    Fase 1 & Espessura mínima de costela    & $t \ge \SI{0.50}{\milli\metre}$                   & NASA-STD-6030 \\
    Fase 1 & Folga livre de despoeiramento  & $\text{gap}_{\text{desp}} \ge \SI{1.50}{\milli\metre}$ & DfAM Standard \\
    \midrule
    Fase 2 & Coeficiente de Poisson efetivo & $\nu_{\text{eff}} < 0{,}0$ (efeito auxético)      & Gibson-Ashby (1997) \\
    \midrule
    Fase 3 & Primeira frequência própria    & $f_1 \ge \SI{100.0}{\hertz}$                      & NASA GEVS-7000A \\
    Fase 3 & Aceleração PSD de entrada      & $\Grmsbase = \SI{14.1}{\Grms}$                    & NASA GEVS-7000A \\
    \midrule
    Fase 4 & Acurácia do metamodelo neural  & $R^2 \ge \num{0.95}$ e $\text{NRMSE} \le \SI{5.0}{\percent}$ & Validação cruzada \\
    Fase 4 & Margem de segurança estrutural & $MS_{\text{yield}} \ge 0{,}0$ ($FS = \num{1.25}$) & NASA-STD-5020B \\
    Fase 4 & Dano de fadiga de Steinberg    & $D \le \num{0.25}$ (fator de vida $4\times$)      & NASA-HDBK-7005 \\
    Fase 4 & Alívio de massa estrutural     & $\Delta m \le \SI{-40.0}{\percent}$ frente ao bloco sólido & CalPoly CDS \\
    \bottomrule
  \end{tabularx}
  \fonte{Elaborada pelo autor (2026).}
\end{table}
